\section{Introduction}
\label{sec:intro}

% @p p-1
Function-as-a-Service (FaaS) platforms bill tenants per invocation and scale
each function down to zero instances when it is idle~\cite{jonas2019cloud}.
The price of scaling to zero is the cold start: the first request after an idle
period waits while the platform boots a sandbox, loads a language runtime and
initializes the function's code. Production traces show that cold starts are
frequent, because most functions are invoked rarely~\cite{shahrad2020serverless}.

% @p p-2
Providers mitigate cold starts with warm pools that keep recently used instances
alive for a fixed keep-alive window. Warm pools trade memory for latency, and the
trade is poor for the long tail of rarely invoked functions: an instance that is
kept warm for ten minutes and then evicted unused costs memory and saves nothing.

% @p p-3
Snapshot restore is the main alternative. The platform checkpoints an initialized
instance once and restores it from the snapshot on every cold start, which skips
runtime and application initialization~\cite{ustiugov2021benchmarking}. Restore
is not free, however: the restored instance faults in its memory one page at a
time, and these page faults dominate restore latency.

% @p p-4
This paper makes two observations. First, the set of pages a function touches is
stable across invocations. Second, that set is small relative to the snapshot.
We build on both in \textsc{Prefetch} (Section~\ref{sec:method}) and evaluate it
in Section~\ref{sec:results}.
